% ============================================================
\section{Synthesis, Falsifiability, and Current Frontier}
\label{sec:discussion}
% ============================================================

\subsection{Temporal Realization and Universal Consequences}

Part~I proves the structural calculus and continuum bridge. Part~II assumes
the explicit three-law physical hypothesis and proves its carrier-level
consequences. Within that boundary, the bedrock counting and boundary layers,
the horizon/tower calculus, the horizon complex, the carrier-forcing
theorems, and the continuum bridge culminate in the temporal interface
exported by the unique \(D=3\) residual. The Part~II consequences on the
faithful selected cut are then derived from temporal realization, terminal
exactification, and endpoint completion.

On the realized cut class, the temporal realization law says that physical
time is the causal realization of ordered persistence on a single observed
bundle. That bundle carries one exact visible law; visible source is only
exchange or returned residual on that same bundle; observable irreversibility
is unrecovered residual; and no exogenous constitutive completion is admitted.
Least motion selects the canonical representative among faithful cuts carrying
the same ordered persistence.

The carrier-level consequences are then read on that common selected bundle.
Within the present scope, the intrinsic laws are exhausted by
the three primitive direct channels together with the unique geometric lift of
the primitive \(3\)-channel. Their recognizable normal forms are the
Schr\"odinger, wave/Klein--Gordon, Maxwell, and Einstein equations.
Hydrodynamic and thermodynamic laws are physical as well, but on the present
scope they are read only as effective closures of unresolved return rather
than as further fundamental laws. This is the closed classification result
under the current Part~II hypotheses.

The hypothesis levels also remain explicit. The exact physical field and its
carrier-specific normal forms are proved from the explicit three-law Part~II
hypothesis fixed in the front claim boundary and Chapter~10. The auxiliary
completion-side normal form is retained only as an alternative presentation on
the same selected bundle. The weaker closure-current / event-exactification
story remains an upstream motivation route, not the formal hypothesis of the
proved Part~II theorems.

\begin{tcolorbox}[colback=white, colframe=black!70, title={Present scope at a glance}]
\begin{itemize}[itemsep=0.2em, leftmargin=1.5em]
\item The narrow bedrock-certified claim concerns only the rebuilt Part~I spine.
\item The continuum bridge and the Part~II theorem hypotheses are separately audited on the Lean PartIII and PartIV developments.
\item The Part~II theorems are not derived from Part~I alone and not from the first two laws alone.
\item On the present scope, the complete fundamental set is the three primitive direct laws together with the unique geometric lift.
\item The recognizable Schr\"odinger, wave/Klein--Gordon, Maxwell, and Einstein equations are the intrinsic realizations of that fundamental set.
\item Hydrodynamic and thermodynamic laws are effective closures on the same selected-cut grammar rather than additional fundamental laws.
\item Visible source on the selected cut is exchange or returned residual, not primitive creation.
\item The explicit Part~II closure theorem now uses the endpoint completion law as the third physical law on top of temporal realization and terminal exactification.
\item Returned-remainder theorems identify exact correction slots around the leading-order laws, but their empirical content depends on fixing a concrete realized model or selected cut.
\item Coupling constants, mass ratios, and the particle spectrum remain open.
\end{itemize}
\end{tcolorbox}

\subsection{Falsifiability, Empirical Stakes, and Structural Predictions}

The present framework makes structural rather than quantitative predictions.
The strongest ones are falsifiable in a clean way.

\begin{enumerate}[label=(\roman*), itemsep=0.3em]
\item The number of primitive direct carrier types is exactly three: phase,
      wave, and gauge. A genuinely fourth primitive carrier, not reducible to
      the intrinsic \(1+2+3\) channel algebra, would falsify the present
      channel-forcing picture.
\item Geometry is predicted to enter as the unique symmetric-response lift of
      the primitive \(3\)-channel rather than as an independent primitive
      sector. A geometric law class shown to be algebraically independent of
      that lift would falsify the present emergence claim.
\item On the selected cut, visible source is predicted to be exchange or
      returned residual, not primitive creation. Observation of genuine
      primitive visible creation on a closed selected sector would falsify the
      exchange-balance clause of the temporal realization hypothesis.
\item Hydrodynamics is predicted to remain an effective closure law, not a
      primitive direct channel. A fundamental fluid law that could not be read
      as resolution-integrated closure would be a counterexample to the
      current scope.
\end{enumerate}

These structural predictions are the cleanest present empirical stakes of the
framework, because they do not require fixing a particular realized model
beyond the selected-cut theorem hypotheses themselves. There is also a second, more
model-dependent predictive layer. The returned-remainder theorems for the
intrinsic equations identify exact correction slots around the leading-order
Schr\"odinger, wave/Klein--Gordon, Maxwell, Einstein, and hydrodynamic forms.
Once a concrete realized cut is fixed, those returned remainders become
genuine correction predictions rather than mere bookkeeping terms. What the
present manuscript does \emph{not} yet claim is a universal numerical
identification of those correction slots across all realizations.

The present scope does not yet include numerical predictions of coupling
constants, mass ratios, or the particle spectrum. The channel dimensions
\(1{+}2{+}3\) are structurally suggestive of the Standard Model gauge
structure~\cite{PeskinSchroeder1995}, and the vacuum topology of the three
channels provides a separate candidate derivation of the fine structure
constant. Another natural open problem is whether the endpoint completion law
is genuinely independent, or instead already forced by temporal realization
and terminal exactification through a sharper intrinsic completion theorem on
the fixed assembly. These remain open problems for the next stage
of development.

Chapter~15 sharpens the same boundary from a different direction: it isolates
the pre-dynamical temporal interface structure on which the temporal realization
hypothesis is read and records that this physical-law layer is portable to any
external host that realizes that same structure. Accordingly, the central
physical-law claim is now stated at
the interface level rather than as a literal requirement that the world first
be accepted as fundamentally discrete.

At the present theorem level, the physical-law completion is also now rigid in a
scoped sense. Proposition~\ref{prop:scoped-temporal-completion-unique} shows
that once a least-motion completion satisfies the temporal realization
principle, no additional independent physical-law data remain on the selected
cut beyond horizon-preserving equivalence, while
Corollary~\ref{cor:exogenous-law-side-alternatives-excluded} records the
corresponding exclusion of primitive visible source, exogenous constitutive
completion, and non-return-based forcing.

\subsection{Intrinsic and Effective Forms}

The lower bedrock of Distinction, Free Composition, and Realized Closure,
together with its derived Constancy / Existence / Incidence / Closure-Balance
structure, the characteristic realization of Chapter~13, the continuum bridge
at the end of Part~I, and the temporal realization hypothesis of Part~II yield one
common characteristic field on the horizon complex and, from it, two kinds of
visible law outcome. On the same selected realization, one common selected
state-field law is read through different carriers. The resulting
Schr\"odinger, wave/Klein--Gordon, Maxwell, and Einstein equations are
therefore presented here as realizations of one common field rather than as
four unrelated symmetry recoveries. On the
least-motion phase carrier the resulting physical law class is a reversible
phase-Hamiltonian law; when its canonical representative is written in the
standard phase representative and Laplace-type transport form, it is the
Schr\"odinger equation. On the least-motion reversible wave carrier the
resulting physical law class is a scalar second-order wave law; in
Laplace-type transport form with scalar mass/potential closure, it is the
scalar wave/Klein--Gordon law. On the least-motion vector/one-form gauge
carrier the resulting physical law class is the curvature/exchange law; when
its selected-cut curvature adjoint is read through Hodge duality
\cite{Maxwell1865,Jackson1999}, it is the Maxwell law. On the least-motion
symmetric-tensor carrier the resulting
physical law class is the covariant geometric law; when its canonical
trace-adjusted curvature representative is read through Levi-Civita/Bianchi
duality \cite{Einstein1915,Wald1984}, it is the Einstein equation. These are
presented as the intrinsic
visible laws. By contrast,
on the least-motion scalar/vector kinetic carrier the exact intrinsic content
is the moment-balance hierarchy, while the recognizable continuity-plus-momentum
sector of compressible Navier--Stokes is presented as the effective
hydrodynamic closure obtained after unresolved kinetic return is compressed
into endpoint response.

The same law set also fixes the admissible meaning of visible sources. Closed
reversible phase and wave sectors are source-free at leading order, so any
visible inhomogeneity there is an exchange remainder induced by coupling or by
hidden-sector return. On the kinetic carrier, moment production is visible
exchange with unresolved higher moments and vanishes on collision-invariant
channels. On the gauge carrier, the one-form \(J\) is a visible exchange
current, and on the geometric carrier the tensor \(T\) is a visible exchange
tensor furnished by a preceding selected law. This is why the geometric law
is not being treated as a constitutive closure in the same sense as
hydrodynamics: the Einstein-side law is geometric response, while any
constitutive closure issue belongs to the matter sector furnishing \(T\), not
to gravity itself. The same effective-closure analysis now also isolates the
mechanically forced thermodynamic content already present in Part~II:
returned-flux / unrecovered-stock balance, monotone irreversibility
bookkeeping, a genuine selected-history second law on the selected coarse
path, and now also an exact selected-tower thermodynamic state space with its
canonical entropy functional, exact entropy/stock decomposition law, and
equilibrium-tail criterion. The remaining thermodynamic extension concerns
broader macrostates, local equilibrium, and broader constitutive closure
beyond the selected refinement history.

\subsection{Closure of the Present Scope}

Each intrinsic law passes through the same closure mechanism: local
candidates are reduced to a finite closure class, compatibility and
symmetry collapse that class, and minimal order selects the surviving
endpoint class. On the reversible intrinsic laws, exact energy-orthogonal
bookkeeping on the realized cut then fixes the canonical trial/test pairing
on the generated carrier and thereby the unique strong representative of that
rigid endpoint class. On the present scope, those intrinsic
endpoint-rigidity theorems are now in place. The phase law is closed by
Theorem~\ref{thm:phase-local-endpoint-collapse}, the wave law is closed to its
Laplace/scalar-profile form by
Theorem~\ref{thm:wave-local-endpoint-collapse}, and the gauge and geometric
laws are closed by
Theorems~\ref{thm:gauge-local-endpoint-collapse} and
\ref{thm:geometric-local-endpoint-collapse}. The optional exact constant-mass
Klein--Gordon refinement is likewise closed on homogeneous selected-wave
sectors by
Theorem~\ref{thm:homogeneous-wave-scalar-profile-forces-constant-mass}.

These closures are local to already forced structural frontiers. The phase
law is already reduced to its canonical complex/Laplace/potential form, the
wave law to its canonical Laplace/scalar-profile form, the gauge
law to its canonical curvature-adjoint form with Maxwell as the standard
Hodge form \cite{Jackson1999}, and the geometric law to its canonical
trace-adjusted curvature form with Einstein as the standard
Levi-Civita/Bianchi form \cite{Wald1984}. By
Theorem~\ref{thm:temporal-realization-forces-selected-cut-forcing-package},
Theorem~\ref{thm:part-iii-realized-horizon-tower-forces-selected-schur-hft2-return-package},
and
Corollary~\ref{cor:temporal-realization-on-part-iii-horizon-tower-forces-grounded-selected-schur-package},
the selected-cut realization theorem and the selected Schur/HFT-2 return theorem
are now closed together on the same least-motion data. Subsequent extensions
concern broader carrier families, broader effective closures, and
thermodynamic development beyond the selected tower. The remaining major
presentation task is not to discover a new law, but to write this same
characteristic-field law in its final external PDE/action language.

\subsection{Bundle Structure and Forcing}

\subsubsection{Selected Bundle}

The common object in Part~II is the selected visible law together with its
inherited endpoint closure class on the least-motion cut. When the
faithful-visibility hypotheses of
Theorem~\ref{thm:common-selected-bundle-assembly} are met, the four intrinsic
carriers together with the hydrodynamic balance carrier become
realized sectors of one selected observed bundle. These carriers are not
imposed by an external dictionary. Rather,
Theorem~\ref{thm:canonical-sector-generation-reduction} identifies them as the
canonical minimal compatible realizations of the intrinsic \(1+2+3\) channel
algebra, together with the canonical symmetric-response lift of the primitive
\(3\)-channel and the separate resolution-integrated balance closure law. A
separate \(D=2\) law also appears within the earlier structural chapters: the unit-balanced seed
carries a one-channel primitive-support tower, maximal cyclic refinement
proceeds by prime-index steps, and the associated spectral defect family has a
unique continuum form \(\mathcal E_{\infty}^{\mathrm{prim}}\). The
corresponding arithmetic bridge \(\mathcal E_{\infty}^{\mathrm{prim}}
= W_{\xi}\) is external to the calculus.

\subsubsection{Intrinsic Channel Structure}

Within the current intrinsic family, the structural ontology is fixed. The
three primitive persistence channels and the geometric lift exist on the
selected observed bundle, are pairwise distinct, and exhaust the intrinsic
family. Each intrinsic carrier therefore realizes exactly one of
those roles. The canonical selected Schur bridge then drives mode selection
internally: Schur-induced candidate classes collapse to singleton reduced
selection data modulo horizon-preserving equivalence, the intrinsic
symmetry/order argument forces the three direct channel classes, and the
primitive \(3\)-channel furnishes the canonical symmetric bilinear lift whose
symmetric-square response closes the geometric law. With
Theorem~\ref{thm:three-forcing-bridges-from-bedrock-to-primitive-physics} and
Corollary~\ref{cor:current-forcing-frontier-for-primitive-physics}, the
intrinsic classification is closed on the present scope:
phase, wave, gauge, and geometry are all forced on one selected bundle, and
the exact constant-mass Klein--Gordon refinement is additionally closed on
homogeneous selected-wave sectors.

Across the intrinsic carrier sections, the governing logic stays the same; what varies
is the visible carrier and its inherited compatibility identity: reversible
phase transport on the phase carrier, reversible second-order transport on the
wave carrier, resolution-integrated balance hierarchy on the kinetic carrier,
nilpotent exactness on the
one-form carrier, and covariant divergence compatibility on the
symmetric-tensor carrier. In each case the same scheme is used: exact
coherent law, temporal realization of persistent exactification on the least-motion
cut, residual return on that cut, physical realization of the selected visible
law, exchange-balanced source admissibility, characteristic endpoint reduction,
canonical sector generation,
symmetry-rigid minimal compatible closure, and finally carrier-level physical
law identification.

\subsubsection{Bundle-First Decomposition}

Part~II therefore yields a common structural description in which the intrinsic
law classes arise as realizations of the primitive \(1+2+3\) channel
structure and its canonical geometric lift. In the bundle-first form of
Theorem~\ref{thm:scoped-bundle-first-coherence-decomposition}, the common
selected bundle itself decomposes, on the present scope, into exactly
those three primitive direct channels together with that geometric lift, while
hydrodynamic/macroscopic laws remain effective closures above them rather than
additional primitive coherence types. The stronger
exact-visible general-bundle statement is now isolated as
Theorem~\ref{thm:general-bundle-exact-visible-exhaustion}: once every
irreducible exact-visible self-sufficient law class on the selected bundle is
shown to be either Schur-induced for one of those four intrinsic roles or a
resolution-integrated effective balance closure, no further primitive visible
law type remains on that scope. The newest reduction sharpens this once
more: by
Theorem~\ref{thm:direct-lift-closure-trichotomy-forces-intrinsic-role-dichotomy}
and
Corollary~\ref{cor:direct-lift-closure-trichotomy-closes-general-bundle-exhaustion},
it is enough to derive the more physical general-bundle direct/lift/closure
trichotomy. On the present scope, canonical realization of the
selected bundle upgrades that trichotomy to the intrinsic-role dichotomy, and
role-compatible exact-visible law classes are automatically Schur-induced.

\subsubsection{Common Theorem Spine}

Three common reductions govern the selected bundle. Theorem~
\ref{thm:common-selected-bundle-assembly} shows that once the four intrinsic
carriers together with the hydrodynamic balance law are faithfully
visible on one physically realized least-motion cut, they are associated
visible carriers of one selected observed bundle and are governed by the same
selected visible effective law. Theorem~
\ref{thm:scoped-intrinsic-sector-generation} proves the non-dictionary bridge
from the intrinsic \(1+2+3\) channel algebra to the four intrinsic visible
sectors: once the intrinsic family is visible on one selected bundle,
those sectors are generated intrinsically there as the unique minimal
compatible realizations of the three primitive persistence channels together
with the canonical symmetric-response lift of the primitive \(3\)-channel.
The hydrodynamic law is then read separately on that same bundle as a
resolution-integrated effective closure rather than as a further primitive
channel or intrinsic lift. Theorem~
\ref{thm:three-forcing-bridges-from-bedrock-to-primitive-physics} now states
the forcing program in one theorem: bedrock and scalar-recursive continuum
forcing determine the temporal interface, temporal realization forces the three
primitive visible channels, and the primitive \(3\)-channel forces geometry as
its unique intrinsic lift. Theorem~
\ref{thm:first-three-intrinsic-law-classes-already-forced} and Theorem~
\ref{thm:three-primitive-channels-already-forced} make the carrier-by-carrier status
explicit: the three primitive channels are already forced on the selected
bundle through the phase, wave, and gauge laws. Theorem~
\ref{thm:canonical-selected-schur-bridge-forces-primitive-three-bilinear-lift-package}
and
\ref{thm:primitive-three-symmetric-bilinear-lift-forces-schur-tensor-lift}
then derive the canonical bilinear self-coupling of the primitive
\(3\)-channel and its symmetric-square lift directly on the selected cut.
Theorem~
\ref{thm:lane-by-lane-intrinsic-closure-frontier} sharpens that separation
again by recording that the intrinsic laws are now all closed: phase,
wave, gauge, and geometry are discharged by
Theorems~\ref{thm:phase-local-endpoint-collapse},
\ref{thm:wave-local-endpoint-collapse},
\ref{thm:gauge-local-endpoint-collapse}, and
\ref{thm:geometric-local-endpoint-collapse}. Theorem~
\ref{thm:geometry-unique-lift-of-primitive-three-channel} records the
ontological conclusion of that step: geometry is not a fourth primitive
channel, but the unique intrinsic lift of the primitive \(3\)-channel. Theorem~
\ref{thm:common-endpoint-rigidity-reduction} reduces the remaining
carrier-level physical problem to identification of a natural representative of
an already forced channel-scalar endpoint closure class. Theorem~
\ref{thm:coupled-exchange-balance} gives the common interaction law on coupled
selected carriers: visible sources cancel pairwise as balanced exchange, and
the two-endpoint case collapses to equal-and-opposite interaction. On the
bridge side, Theorem~\ref{thm:explicit-continuum-verification-package}
collects the existing explicit dyadic, local-ray, fibered-promotion, and
primitive-support model families, furnishing explicit continuum verification.
Theorem~\ref{thm:part-iv-law-completion} then states the exact criterion for a
completed law chapter: complete the canonical selected
Schur bridge, prove the common selected bundle intrinsically, generate the
sectors intrinsically from the \(1+2+3\) channel algebra, and complete the
carrier-level natural-operator identifications.

Theorem~\ref{thm:defect-tower-complete-carrier-forcing} then assembles the
chain: the \(S_4\)-irreducible tower on \(V_{\mathrm{slot}}\) is read as a
three-step horizon tower, and \(\partial^2=0\) forces the four carrier types
through the defect mechanism, the HFT-2 energy budget, and the scalarization
theorems already available in the structural calculus.
Corollary~\ref{cor:single-source} records the single-source identification:
the common carrier taxonomy and closure grammar derive from the single
structural identity \(\partial^2=0\) on the canonical balanced clique complex,
while the intrinsic equations appear later as autonomous normal forms under the
explicit endpoint-rigidity and representative hypotheses.

\subsection{Dependency Audit of the Intrinsic Equations}

The intrinsic equations are clearest when their universal core is separated from
their carrier-identification data. Table~\ref{tab:flagship-dependency-audit}
records that split on the present scope. The first column after the equation family
contains the common law-side content already supplied by the preceding theorem
chain. The next column contains only carrier-identification data. The final
column is the representative form. Accordingly, the intrinsic equations
are one common temporal structure with several carrier realizations, not several
unrelated dynamical laws.

\begin{table}[ht]
\centering
\scriptsize
\renewcommand{\arraystretch}{1.15}
\setlength{\tabcolsep}{3pt}
\begin{tabular}{@{}P{1.7cm}P{4.35cm}P{4.0cm}P{4.5cm}@{}}
\toprule
\textbf{Equation family} & \textbf{Universal core already supplied by the preceding theorem chain} & \textbf{Carrier-identification data} & \textbf{Representative form} \\
\midrule
Phase
& selected cut, selected visible effective law, common selected bundle,
  primitive \(1\)-channel forcing, exchange/return grammar, and endpoint-rigidity
  schema
& least-motion reversible phase carrier, energy-orthogonal phase pairing, and
  phase-compatible local transport family
& canonical complex/Laplace/potential representative; Schr\"odinger form plus
  returned phase remainder \\

Wave
& same selected cut and bundle, primitive \(2\)-channel forcing, common
  reduction to a rigid endpoint class, and same return grammar
& reversible second-order scalar carrier, no-preferred-direction compatibility,
  and pairing-symmetric transport family
& Laplace/scalar-profile wave law, Klein--Gordon on homogeneous sectors, and
  returned wave remainder \\

Gauge
& same selected cut and bundle, primitive \(3\)-channel forcing, common
  exchange/return grammar, and same endpoint-rigidity machinery
& one-form comparative carrier, nilpotent exactness, and gauge-compatible
  minimal-order response
& canonical curvature-adjoint representative; Maxwell form plus returned
  gauge current \\

Geometry
& same selected cut and bundle, geometric lift of primitive \(3\), common
  law-side closure grammar, and same endpoint-rigidity reduction
& faithful symmetric bilinear lift, divergence-compatible symmetric response,
  and symmetric-tensor carrier
& canonical trace-adjusted curvature representative; Einstein form\newline
  and current-solved remainder model linearized on a transverse-traceless
  sector \\

Hydro\-dynamics
& same selected data, selected visible law, residual-return grammar, and
  effective-closure reading of unresolved sectors
& scalar/vector kinetic carrier, moment-balance hierarchy, and isotropic
  first-order closure form
& continuity/momentum closure with compressible
  Navier--Stokes-type stress law\newline plus exact hydrodynamic
  remainder \\

Thermo\-dynamics
& same selected tower, returned flux / unrecovered stock bookkeeping,
  irreversibility grammar, and selected-history ordering
& selected coarse history / selected-tower state space together with the
  exact entropy-stock decomposition
& selected-history second law, canonical selected-tower entropy,
  and exact equilibrium-tail criterion \\
\bottomrule
\end{tabular}
\caption{Dependency audit for the intrinsic equations. The universal dynamical
core is common across the equation families; carrier-specific material
identifies the visible carrier on which that core is realized, and the final
column records the canonical representative form.}
\label{tab:flagship-dependency-audit}
\end{table}

\begin{remark}[What the dependency audit shows]\lean{flagship_law_completion}
Table~\ref{tab:flagship-dependency-audit} gives the clearest answer to the
question of whether the intrinsic equations import hidden extra dynamics. On the
present scope they do not. What varies from one equation family to another is the visible carrier
and its compatibility identity, not the law-side completion principle itself.
The remaining local assumptions are carrier-identification and endpoint
realization assumptions.
\end{remark}

\subsection{Wave Example}

The reversible wave carrier yields an explicit observational consequence. In
the nearest-neighbor massless realization of
Section~\ref{sec:flagship-wave}, plane-wave modes satisfy
\[
\omega(k)^2=\frac{4c_\chi^2}{\ell_\chi^2}\sin^2\!\left(\frac{k\ell_\chi}{2}\right),
\]
so that with \(E=\hbar\omega\) and \(E_\chi:=\hbar c_\chi/\ell_\chi\),
\[
v_g(E)=c_\chi\sqrt{1-\left(\frac{E}{2E_\chi}\right)^2}
\qquad (0\le E\le 2E_\chi).
\]
In the low-energy regime this yields the universal asymptotic form
\[
\frac{\delta v}{c_\chi}\approx -\frac18\left(\frac{E}{E_\chi}\right)^2.
\]
The sign is fixed: high-energy components are subluminal. The same wave law
also predicts a hard cutoff at \(E=2E_\chi\), where the group velocity
collapses to zero.

That lossless nearest-neighbor signature is now accompanied by an exact
non-lossless refinement. By
Theorem~\ref{thm:exact-schur-corrected-wave-dispersion}, whenever the selected
wave remainder is mode-compatible on a plane-wave sector, the observed
dispersion law becomes
\[
\omega(k;z)^2
=
m_{\chi}^2
+
\frac{4c_{\chi}^2}{\ell_{\chi}^2}
\sin^2\!\left(\frac{k\ell_{\chi}}{2}\right)
-
\sigma_{\chi}(k;z),
\]
with the corresponding exact group-velocity correction determined by the same
returned wave profile \(\sigma_{\chi}(k;z)\). The clean lattice law is
therefore the special case \(\sigma_{\chi}(k;z)=0\), not a separate equation.
Theorem~\ref{thm:hidden-channel-wave-remainder-and-branch-splitting} then gives
an exact solved model of that correction structure: a locally coupled
hidden wave channel produces the explicit self-energy
\[
\sigma_{\chi}^{\mathrm{hid}}(k,\omega)
=
\frac{\gamma_{\chi}^2}{\omega_{\mathrm{hid}}(k)^2-\omega^2},
\]
and the visible wave law reorganizes into two exact mixed branches
\(\omega_{\pm}(k)\). So for the wave equation the returned remainder is already
more than a symbolic slot; it can be computed exactly on a concrete realized
family. Corollaries~\ref{cor:exact-off-resonant-hidden-channel-branch-shifts},
\ref{cor:exact-avoided-crossing-coupled-wave-branches}, and
\ref{cor:low-energy-renormalized-mass-and-wave-speed-hidden-channel} then show
exactly what that solved remainder does: it shifts the visible branch away
from resonance, opens an avoided crossing of size \(2\gamma_{\chi}\) in
squared-frequency variables at resonance, and renormalizes the visible mass
gap and long-wave propagation speed at low energy.

\subsection{Returned-Remainder Program}

The present framework is not exhausted by recovery of recognizable lossless
laws. Theorem~\ref{thm:returned-remainders-form-predictive-correction-surface}
shows that, on the selected bundle, the textbook-looking standard identities
are precisely the leading-order representatives obtained when the exact
returned remainders vanish. The nonzero remainder terms are therefore not
arbitrary error bars or phenomenological add-ons. They are the native
correction structure of the theory: exact returned-residual data induced by the
same coherent bundle and the same selected cut.

On the present scope this means that the genuinely predictive work lies in
computing the induced scales and returned remainders on concrete realized cuts:
\(\mathsf R_{\mathrm{ph}}\) for the phase equation,
\(\mathsf R_{\mathrm{wav}}\) for the wave equation,
\(\mathsf R_{\mathrm g}\) for the gauge equation,
\(\mathsf N_{\mathrm{grav}}\) for the geometric equation, and
\(\mathsf R_{\mathrm{hyd}}\) for the resolution-integrated hydrodynamic closure.
These terms measure the exact visible effect of unresolved return; they are the
place where the theory can predict departures from clean Schr\"odinger, wave /
Klein--Gordon, Maxwell, Einstein, and hydrodynamic normal forms.
Equivalently, the intrinsic equations now carry explicit symbolic correction terms
\[
\Sigma_{h,\mathrm{ph}}(z),\quad
\Sigma_{h,\mathrm{wav}}(z),\quad
\Sigma_{h,\mathrm g}(z),\quad
\Sigma_{h,\mathrm{grav}}(z),\quad
\Sigma_{h,\mathrm{hyd}}(z),
\]
namely the carrier-level images of the one common selected Schur correction
\(\CCSigma{h}{z}\) under the corresponding endpoint identifications.

\begin{figure}[ht]
\centering
\resizebox{\textwidth}{!}{%
\begin{tikzpicture}[x=1cm,y=1cm]
  \node[ccbox, align=center, minimum width=5.2cm, minimum height=0.95cm] (schur) at (0,4.05)
    {\textbf{One selected Schur correction}\\\(\CCSigma{h}{z}\)};
  \node[ccbox, align=center, minimum width=5.8cm, minimum height=0.95cm] (id) at (0,2.7)
    {\textbf{Carrier-level endpoint identifications}\\map \(\CCSigma{h}{z}\) to equation-specific corrections};

  \draw[ccarrow, line width=0.5pt] (schur.south) -- (id.north);

  \node[ccbox, align=center, text width=2.85cm, minimum height=1.25cm] (ph) at (-4.2,1.0)
    {\textbf{Phase}\\Schr\"odinger-type law\\\(+\,\Sigma_{h,\mathrm{ph}}(z)\)};
  \node[ccbox, align=center, text width=2.85cm, minimum height=1.25cm] (wav) at (0,1.0)
    {\textbf{Wave}\\wave / KG law\\\(+\,\Sigma_{h,\mathrm{wav}}(z)\)};
  \node[ccbox, align=center, text width=2.85cm, minimum height=1.25cm] (gau) at (4.2,1.0)
    {\textbf{Gauge}\\Maxwell law\\\(+\,\Sigma_{h,\mathrm g}(z)\)};
  \node[ccbox, align=center, text width=3.05cm, minimum height=1.25cm] (geo) at (-2.2,-0.95)
    {\textbf{Geometry}\\Einstein law\\\(+\,\Sigma_{h,\mathrm{grav}}(z)\)};
  \node[ccbox, align=center, text width=3.05cm, minimum height=1.25cm] (hyd) at (2.2,-0.95)
    {\textbf{Hydrodynamics}\\NS-type closure\\\(+\,\Sigma_{h,\mathrm{hyd}}(z)\)};

  \draw[ccarrow, line width=0.5pt] (id.south) -- ++(0,-0.18) -| (ph.north);
  \draw[ccarrow, line width=0.5pt] (id.south) -- ++(0,-0.18) -- (wav.north);
  \draw[ccarrow, line width=0.5pt] (id.south) -- ++(0,-0.18) -| (gau.north);
  \draw[ccarrow, line width=0.5pt] (id.south) -- ++(0,-0.18) -| (geo.north);
  \draw[ccarrow, line width=0.5pt] (id.south) -- ++(0,-0.18) -| (hyd.north);
\end{tikzpicture}%
}
\caption{The common predictive structure of Part~II. One selected Schur
correction \(\CCSigma{h}{z}\) is carried through the endpoint identifications
to produce the equation-specific returned remainders. The recognizable textbook
laws are the lossless leading-order representatives obtained when those
corrections vanish.}
\label{fig:returned-remainder-surface}
\end{figure}

Four of those intrinsic equations already admit exact solved coherent models on
the present scope. For the wave equation,
Theorem~\ref{thm:hidden-channel-wave-remainder-and-branch-splitting} and its
corollaries compute the returned profile explicitly and identify exact branch
shifts, avoided crossing, and low-energy renormalization. For the phase equation,
Theorem~\ref{thm:hidden-channel-phase-remainder-and-spectral-splitting} and
Corollaries~\ref{cor:exact-off-resonant-phase-line-shifts},
\ref{cor:exact-avoided-crossing-coupled-phase-branches}, and
\ref{cor:low-energy-renormalized-phase-offset-and-transport-coefficient}
likewise compute the coherent returned phase profile on a two-channel model
and show exact spectral line shifts, branch repulsion, and low-energy
renormalization of the visible phase offset and transport coefficient. For the
gauge equation, Theorem~\ref{thm:hidden-channel-gauge-current-and-branch-splitting}
and Corollaries~\ref{cor:exact-off-resonant-gauge-pole-shifts},
\ref{cor:exact-avoided-crossing-coupled-gauge-branches}, and
\ref{cor:low-frequency-induced-exchange-susceptibility} compute the returned
gauge current on a two-channel model and show exact gauge-pole shifts, branch
repulsion, and low-frequency induced exchange susceptibility on the visible
gauge sector. For the geometric equation,
Theorem~\ref{thm:hidden-channel-geometric-remainder-and-branch-splitting} and
Corollaries~\ref{cor:exact-off-resonant-geometric-pole-shifts},
\ref{cor:exact-avoided-crossing-coupled-geometric-branches}, and
\ref{cor:low-frequency-induced-geometric-susceptibility} compute the returned
linearized geometric profile on a two-channel model and show exact
geometric-pole shifts, branch repulsion, and low-frequency induced
symmetric-response susceptibility on the visible geometric sector.

The theory does not identify a nonzero remainder with any named phenomenon in
advance. It does not, by itself, say
that a given correction is ``dark matter,'' ``quantum gravity,''
``decoherence,'' or ``turbulence.'' Those are concrete realization problems.
What the theory does say is sharper: on the present scope, any such
departure must enter through exact returned remainders or induced coefficient
scales on the same selected realization, not through ad hoc primitive source
terms foreign to the bundle. The explicit wave-speed correction and hard cutoff
above are the first concrete example of that predictive structure rather than the
last.

\subsection{Candidate Phenomenology of Returned Remainders}

The theorems above do not identify any nonzero returned remainder with a named
physical anomaly. They do provide a disciplined
phenomenology map: if a concrete realized family yields a nonzero equation-specific
correction \(\Sigma_{h,\bullet}(z)\), then the resulting observable departure
from the corresponding leading textbook law should fall into a restricted class
of signatures. Table~\ref{tab:returned-remainder-phenomenology} records that
map at the present level of rigor. It is intended as a hypothesis table for
future calculation and comparison, not as a proved identification theorem.

\begin{table}[ht]
\centering
\scriptsize
\setlength{\tabcolsep}{3pt}
\begin{tabular}{@{}P{1.6cm}P{2.7cm}P{4.7cm}P{4.5cm}@{}}
\toprule
Equation family & Correction & Observable signature & Comparison language \\
\midrule
Phase
& \(\Sigma_{h,\mathrm{ph}}(z)\)
& phase drift, line shifts, loss of clean interference, or effective
  environment-coupled evolution on the visible cut
& open-system correction, decoherence-type term, level-shift / self-energy
  correction \\

Wave
& \(\Sigma_{h,\mathrm{wav}}(z)\)
& nontrivial dispersion, nonlocal propagation, high-frequency slowdown, or
  cutoff behavior beyond the clean wave/Klein--Gordon form
& dispersive-medium correction, lattice/discrete propagation correction,
  effective wave self-energy; the nearest-neighbor specialization already gives
  a concrete subluminal high-energy correction and hard cutoff \\

Gauge
& \(\Sigma_{h,\mathrm g}(z)\)
& effective mismatch between visible current and clean Maxwell response,
  polarization-like medium effects, or unresolved exchange-current correction
& exchange-current correction, polarization / medium response,\newline
  vacuum-polarization-type effective term \\

Geometry
& \(\Sigma_{h,\mathrm{grav}}(z)\)
& additional effective stress-response on top of clean Einstein behavior,
  backreaction-like deviation, or visible source mismatch not captured by the
  leading geometric carrier
& effective stress-tensor correction, geometric backreaction,\newline
  dark-sector phenomenology at the level of the visible response law \\

Hydro\-dynamics
& \(\Sigma_{h,\mathrm{hyd}}(z)\)
& stress not captured by Newtonian viscosity, including memory, nonlocality,
  anisotropic closure, turbulent transfer, rarefied-gas correction, or
  non-Newtonian response
& Reynolds stress, LES subgrid stress,\newline Burnett / kinetic-memory
  correction,\newline non-Newtonian constitutive correction \\
\bottomrule
\end{tabular}
\caption{Hypothesis map for the returned-remainder structure. Each row is
conditional: if the corresponding selected Schur correction is nonzero on a
concrete realized family, the observable departure from the leading textbook
law should appear in the indicated signature class.}
\label{tab:returned-remainder-phenomenology}
\end{table}

The table is intentionally one level weaker than the proved intrinsic theorems.
It does not assert that any existing anomaly is already explained. It records
the restricted comparison classes suggested by the theory: compute the relevant
\(\Sigma_{h,\bullet}(z)\) on a concrete selected realization, then compare that
computed correction with the phenomenology in the corresponding row.

\subsection{Extensions}

\begin{enumerate}[leftmargin=1.2em, itemsep=0.3em]
\item \textbf{Further carrier classes and broader thermodynamic extension.} Additional
      mixed and many-body carriers belong naturally in the same
      physical-law-selection format. The present thermodynamic extension is now
      stronger than a first bookkeeping layer: returned flux, unrecovered
      stock, monotone irreversibility, selected-history entropy production, the
      selected-tower thermodynamic state space, its canonical entropy
      functional, and its exact equilibrium class are all mechanically forced
      on the selected cut. The broader thermodynamic program beyond that
      selected tower concerns extension to coarse-state spaces not
      canonically induced by the realized tower, local equilibrium and
      equilibrium manifolds there, and broader
      thermodynamic/macroscopic closure principles under explicit
      coarse-graining hypotheses.
      On the spectral side, the unit-balanced \(D=2\) seed defines a
      one-fiber reversible law with a single-channel monodromy/recursion
      stage, a primitive-support tower with prime-index maximal cyclic
      refinement, and a forced continuum spectral defect form
      \(\mathcal E_{\infty}^{\mathrm{prim}}\). Arithmetic realization would
      enter through the external identification theorem
      \(\mathcal E_{\infty}^{\mathrm{prim}} = W_{\xi}\).
\end{enumerate}

On the present scope, the remaining manuscript work is now largely
downstream rather than conceptual: extend the current carrier family, develop
broader effective closures, and compute quantitative remainder laws for the
intrinsic carriers within the same residual-return formalism.
